\section{Task and training}
\label{sec:task}

\subsection{The cued change-detection environment ($T=29$)}
The agent performs a spatially cued orientation-change-detection task over an episode of $T=29$ frames. Each $50\times50$ frame is divided into the same $10\times10$ grid of $100$ patch tokens the encoder reads. On an early frame a coloured cue indicates one side of the display and is then removed. The cue carries two pieces of information. Its \emph{reliability}---the validity, drawn from $\{0.25,0.50,0.75,1.00\}$---sets the probability that an upcoming orientation change occurs at the cued location; its \emph{colour} sets the reward magnitude available on that trial, so the cue specifies both where to attend and how much the trial is worth. Half of all trials are change trials. On a change trial the orientation of the target patch shifts by a magnitude of up to roughly $64^{\circ}$, with the change onset jittered across frames (the change index falls in roughly $[11,25]$) and the new orientation persisting to the end of the episode; the remaining trials contain no change. Frame-to-frame Gaussian orientation noise ($\sigma=5$) sets the discrimination difficulty, so that small changes must be detected against ongoing sensory fluctuation. On every frame the agent emits a binary action, \emph{wait} or \emph{declare change}. Declaring at or after the true change onset is a hit; declaring before it is a false alarm; waiting out a no-change trial to the end is a correct rejection. The long horizon, the jittered onset, the removed-then-remembered cue, and the graded reliability together make this a test of maintained covert spatial attention rather than of a single-frame discrimination, and the cue-colour reward structure adds a value dimension on top of the spatial one.

\subsection{The shared training harness}
All three configurations are trained by an identical offline actor--critic reinforcement-learning procedure, with the same hyperparameters and the same number of updates; the configurations differ only in the encoder wiring of Section~\ref{sec:architectures}. The actor and a distributional critic are one-dimensional-convolutional heads over the token axis: the actor emits the wait/declare logit pair and the critic estimates a five-quantile distribution over returns. The actor loss is a maximum-a-posteriori policy-improvement (MPO/MAP) term blended with a behaviour-cloning term; the critic is trained with a quantile (distributional) Bellman loss against a periodically hard-copied target network that supplies both the critic's bootstrap target and the policy-improvement reference. Whole episodes are stored in a prioritised replay buffer, with priorities set by per-episode quantile error, and a brief forced-wait warm-up trains the critic on full trajectories before the actor begins to act, so the policy is first improved against an already-informative value estimate. Because the harness, the parameter budget and the update count are held fixed, any difference in outcome across the three configurations is attributable to the single structural property that distinguishes them.
